\scanpage{305}
\mspage{255}
\diagram{\(\displaystyle
i_!\left(\begin{tikzcd}[ampersand replacement=\&,row sep=small,column sep=small]
X_1\arrow[d]\&\\S\&X_2\arrow[l]
\end{tikzcd}\right)=
\left(\begin{tikzcd}[ampersand replacement=\&,row sep=small,column sep=small]
X_1\arrow[d]\&\unsure{\varphi|\varnothing}\arrow[l]\arrow[d]\\S\&X_2\arrow[l]
\end{tikzcd}\right),\)}
la loi fonctorielle étant \unsure{évidente|ici évidente}.

Ce n'est pas encore la fin du \unsure{sorite}\,! Il faut \inserted{encore} interpréter le foncteur composé
\[
\mathrm{Ho}_W(\Psi)\mathrel{\overset{i_*}{\unsure{-|\to}}}\mathrm{Ho}_W(\Psi^{\sim})\xrightarrow{\deleted{\unsure{i}}\varepsilon_3^*}\mathrm{Ho}_W(e)=\mathrm{Ho}_W,
\]
où
\diagram{\begin{tikzcd}[ampersand replacement=\&,row sep=small,column sep=small]
\&\&1\arrow[dl,leftrightarrow]\&\\
{\varepsilon_3\colon \unsure{e}\simeq\{3\}}\arrow[r]\&3\&\&0\arrow[ul]\arrow[dl]\\
\&\&2\arrow[ul]\&
\end{tikzcd}}
est l'inclusion du sommet 3 dans $\Psi^{\sim}$.

Si $\mathrm{Ho}_W(\Psi)$ est interprété, conformément au lemme 1, comme $\mathscr F W_{\mathscr F}^{-1}$, ce foncteur composé n'est autre que
\diagram{\begin{tikzcd}[ampersand replacement=\&,row sep=small,column sep=small]
X_1\arrow[d,dash]\&\&{\longmapsto}\&X_1\times_S X_2\\
S\&X_2\arrow[l]\&\&
\end{tikzcd}}
\diagram{\(\displaystyle\biggl(\text{pour }\left(\begin{tikzcd}[ampersand replacement=\&,row sep=small,column sep=small]
X_1\arrow[d,"{p_1}"']\&\\S\&X_2\arrow[l,"{p_2}"]
\end{tikzcd}\right)\in\mathrm{Ob}\,\mathscr F,\quad\text{i.e. }p_1,p_2\in\mathrm{Fib}_W\biggr).\)}

On veut interpréter ce foncteur comme une $\varprojlim$ $W$-homotopique, \deleted{plus \unsure{précisé|précisément}} \inserted{i.e. \inserted{ici} comme} un \deleted{\unsure{produit|foncteur}} ``produit fibré $W$-homotopique''. De façon précise, \deleted{\unsure{on}} \unsure{c'est-|ceci|ce qui}\unsure{=|-}
% archive sheet number: 305
